\subsection{From $Q_8$ to $2I$: Strengthening the Spinorial Hierarchy}
  \label{subsec:q8-to-2i}

  The admissibility hierarchy observed in the quaternion example is not an
  isolated feature of the smallest non-abelian group.
  The same mechanism persists and becomes more pronounced in larger binary
  groups.

  A natural next step is provided by the binary icosahedral group $2I$, which is
  the double cover of the rotational symmetry group of the icosahedron.
  This group plays a distinguished role among finite subgroups of $SU(2)$ and
  contains $Q_8$ as a subgroup.

  As in the previous section, we consider the undirected Cayley graph associated
  with a symmetric generating set $S$.
  The relational Laplacian of this graph again decomposes into sectors labelled
  by the irreducible representations of $2I$.

  The spectrum of $2I$ contains representation sectors of several dimensions.
  Applying the admissibility envelope derived in
  Section~\ref{subsec:admissibility-bounded-flux},
  \[
    A_\rho^{\max}=\frac{c_\chi}{\sqrt{\lambda_\rho}},
  \]
  one finds that the hierarchy between sectors becomes more pronounced than in
  the quaternion case.

  In particular, the admissible amplitude associated with the lowest
  two-dimensional representation exceeds that of the lowest one-dimensional
  sector by a factor
  \[
    \sqrt{\frac{14}{9}} .
  \]

  This strengthening reflects the richer relational structure of the
  binary icosahedral group.
  While several representation sectors remain kinematically available,
  the bounded-flux constraint dynamically selects those with the smallest
  Laplacian eigenvalues.

  The resulting hierarchy naturally organises the representation sectors into a
  tower of admissible modes whose amplitudes decrease with increasing spectral
  frequency.

  From the viewpoint of relational dynamics, the binary icosahedral group thus
  provides a larger finite structure in which the admissibility mechanism
  remains fully operative.

  Moreover, $2I$ is embedded in the continuous group $SU(2)$.
  This embedding suggests that the hierarchy observed in finite relational
  structures may persist along sequences of increasingly refined graphs
  approaching a continuous limit.

  Before exploring this limit, however, it is useful to examine more closely the
  structural conditions that allow generators to participate in admissible
  configurations.
  These conditions can be expressed in terms of representation characters and
  lead to the notion of neutral generators introduced in the next section.
